\documentclass[10pt,a4paper]{amsart}
\usepackage[margin=19mm]{geometry}
\usepackage{amsmath,amssymb,mathtools}
\usepackage[scr=rsfs,cal=euler]{mathalfa}
\usepackage{xfrac}
\usepackage[unicode,hidelinks,bookmarks=false]{hyperref}
\newcommand{\ie}{i.e.\ }
\newcommand{\cf}{cf.\ }
\newcommand{\resp}{respectively\ }
\newcommand{\Subsection}{\subsection}
\providecommand{\nobreakdash}{\nobreak\hbox{-}}
\setlength{\emergencystretch}{2em}
\multlinegap0pt
\usepackage{bm}
\usepackage{bbm}
\usepackage{dsfont}
\usepackage{xspace}
\usepackage{cancel}
\usepackage{mathtools}
\usepackage{amsthm}
\makeatletter
\@ifundefined{theorem}{\newtheorem{theorem}{Theorem}}{}
\@ifundefined{lemma}{\newtheorem{lemma}[theorem]{Lemma}}{}
\makeatother
\def\pp{{\frak p}}
\title[Symbolic powers via extension]{Symbolic powers via extension}
\author{Sankhaneel Bisui, Haoxi Hu}
\date{Source: arXiv:2410.22581v3 (CC BY 4.0)}
\begin{document}
\maketitle

\noindent\emph{Source and licence.} Symbolic powers via extension. Version arXiv:2410.22581v3 (CC BY 4.0), distributed under the Creative Commons Attribution 4.0 International licence, as recorded in the arXiv licence metadata for this exact version and shown on the abstract page https://arxiv.org/abs/2410.22581v3. Full rights notice: licenses/arxiv-2410.22581-CC-BY-4.0.txt.\par\smallskip

\noindent\textit{Editorial scope.} Counted unit: the Lemma whose source label is \texttt{lemma: tensor of localization} in the article (source0.tex lines 245--247, source label \texttt{lemma: tensor of localization}), delivered together with the article's own complete original proof of it (source0.tex lines 249--275, one \texttt{proof} environment of 27 lines). No mathematics is written, repaired or re-derived: statement and proof are the article's own text line for line. Required context retained uncounted: lemma: extension of an ideal (lines 231--233). Every definition, display and label of those lines is retained unchanged. Omitted from this package and not redistributed: 1--230, 234--244, 248--248, 276--561. Nothing inside a retained excerpt is omitted or altered: every retained body is delivered exactly as the article's lines are written, so no omission receipt of any kind accompanies this package. Citations: the retained excerpts carry no citation command at all, so no bibliography entry of the article is transcribed and no reference is redistributed. Reference adaptation: none was needed and none was made; every reference inside the delivered unit resolves inside it, so no unresolved reference or citation is produced. Printed slips kept literal: no printed word, symbol, hypothesis or number of the counted statement or of its proof was altered. Layout and class: the article's own class is \texttt{amsart} with packages including hyperref, amssymb, verbatim, amscd, amsmath, graphicx, nccmath, enumerate, caption, subcaption, pdfsync, color, colortbl, fullpage; the wrapper is the lane's pinned \texttt{amsart} editorial wrapper at 10pt on A4, which restates the theorem-like environments the retained text needs and adds the article's own macro definitions that the retained text needs (\texttt{\textbackslash pp}), with extra wrapper packages (bm, bbm, dsfont, xspace, cancel, mathtools). The delivered numbering is the unit's own. Assets: no retained span contains a figure environment, a graphics-inclusion command, a PDF-page inclusion, a table or a TikZ picture; no image file is copied into this package, the package declares no asset dependency and no image file is redistributed with it. Licence: version arXiv:2410.22581v3 of this article is distributed under the Creative Commons Attribution 4.0 International licence, as recorded in the arXiv licence metadata for this exact version and shown on the abstract page https://arxiv.org/abs/2410.22581v3. A full rights notice is delivered at licenses/arxiv-2410.22581-CC-BY-4.0.txt. Characters: every retained excerpt is pure ASCII, so the delivered engine, XeLaTeX with its default Latin Modern text font, reports no missing character.
	\begin{lemma}\label{lemma: extension of an ideal}
		Let $M$ be a $R$-module where $R$ is a commutative ring, and $I$ be an ideal of $R$, then $(R/I) \otimes_R M \cong M/IM$. Furthermore, if $M$ is flat over $A$, we also have $I \otimes_R M \cong IM$.
	\end{lemma}

	\begin{lemma}\label{lemma: tensor of localization}
	If a flat ring homomorphism $\phi: A\to B$ extends a prime $\mathfrak{p}$ to a prime $\mathfrak{p}B$, and $\mathfrak{p}$ is a contraction of some ideals in $B$, then $B_{\mathfrak{p}B} \cong A_{\mathfrak{p}} \otimes_A B_{\mathfrak{p}B}$ as $A$-algebra.
	\end{lemma}

	\begin{proof}
	Applying Lemma \ref{lemma: extension of an ideal} to $B_{\mathfrak{p}B}$, we have $B_{\mathfrak{p}B} \cong ( (A \otimes_A B) \backslash (\mathfrak{p}B \otimes_A B) )^{-1} (A \otimes_A B)$. Our goal is to show $( (A \otimes_A B) \backslash (\mathfrak{p}B \otimes_A B) )^{-1} (A \otimes_A B) \cong A_{\mathfrak{p}} \otimes_A B_{\mathfrak{p}B}$.
	
	Define $$\psi_1 : ( (A \otimes_A B) \backslash (\mathfrak{p}B \otimes_A B) )^{-1} (A \otimes_A B) \rightarrow A_{\mathfrak{p}} \otimes_A B_{\mathfrak{p}B} \text{ by } \psi_1\left(\frac{a \otimes b}{c \otimes d}\right) = \frac{a}{c} \otimes \frac{b}{d}.   $$ %where $a \otimes b \in A \otimes_A B$ and $c \otimes d \in (A \otimes_A B) \backslash (\mathfrak{p}B \otimes_A B)$.
	
	Similarly define  $$ \psi_2 : A_{\mathfrak{p}} \otimes_A B_{\mathfrak{p}B} \rightarrow ( (A \otimes_A B) \backslash (\mathfrak{p}B \otimes_A B) )^{-1} (A \otimes_A B) \text{ by } \psi_2\left(\frac{a}{c} \otimes \frac{b}{d}\right) = \frac{a \otimes b}{c \otimes d}.$$ %where $\frac{a}{c} \in A_{\mathfrak{p}}$ and $\frac{b}{d} \in B_{\mathfrak{p}B}$.
	
	First we need to show both maps are well-defined. 
	
	Since $c \otimes d \notin \mathfrak{p} \otimes_A B$, then $c \notin \mathfrak{p}$. We also can consider $c \otimes d = 1 \otimes cd \notin 1 \otimes_A \mathfrak{p}B$, so $cd \notin \mathfrak{p}B$. We also know that $c \notin \mathfrak{p}$, this forces $d \notin \mathfrak{p}B$ since $\mathfrak{p}B$ is prime. Therefore $\psi_1$ is well-defined.
	
 We prove that	$\psi_2$ is also well-defined. Indeed, $c \notin \mathfrak{p}$ and $d \notin \mathfrak{p}B$, then $\phi(c)d \notin \mathfrak{p}B$ because $\phi(A) \cap \mathfrak{p}B = \phi(\mathfrak{p})$. To see this, let $\mathfrak{p} = \mathfrak{q}\cap A$ for an ideal $\mathfrak{q}$ in $B$. From this we  observe that $\mathfrak{p}B\cap A  = \left( \left( \mathfrak{q}\cap A\right) B\right) \cap A  = \mathfrak{q}\cap A=\pp$. This implies that $\phi(A) \cap \mathfrak{p}B = \phi(\mathfrak{p})$. Note that $\phi(A) \cap \mathfrak{p}B \supset \phi(\mathfrak{p})$. Let $y \in \phi(A)\cap \pp B$. Then $y=\phi(x)$ for some $x\in A$. We show that $x \in \pp$. Now since $y \in \pp B$ then $x\in \pp B\cap A=\pp$. 
 
Notice that $\psi_1$ and $\psi_2$ are inverses to each other. Thereofre we just need to show that $\psi_1$ is an $A$-algebra homomorphism, then we are done. Verifying the map commutes with the multiplication is straightforward, we only verify the addition here. 
   Let $\frac{a_1 \otimes b_1}{c_1 \otimes d_1}$ and $\frac{a_2 \otimes b_2}{c_2 \otimes d_2}$ be elements in $( (A \otimes_A B) \backslash (\mathfrak{p}B \otimes_A B) )^{-1} (A \otimes_A B)$. Then
	\begin{align*}
	\psi_1 \left(\frac{a_1 \otimes b_1}{c_1 \otimes d_1}\right) +  \psi_1 \left(\frac{a_2 \otimes b_2}{c_2 \otimes d_2}\right)&=\frac{a_1}{c_1} \otimes \frac{b_1}{d_1} + \frac{a_2}{c_2} \otimes \frac{b_2}{d_2}\\
	&= \frac{a_1 c_2}{c_1 c_2} \otimes \frac{b_1 d_2}{d_1 d_2} + \frac{c_1 a_2}{c_1 c_2} \otimes \frac{d_1 b_2}{d_1 d_2} \\
	&=\left(\frac{1}{c_1 c_2} \otimes \frac{1}{d_1 d_2}\right) \cdot ( a_1 c_2 \otimes b_1 d_2 + c_1 a_2 \otimes d_1 b_2 ).
	\end{align*}
Again, 	
\begin{align*}
\psi_1 \left( \frac{a_1 \otimes b_1}{c_1 \otimes d_1} + \frac{a_2 \otimes b_2}{c_2 \otimes d_2} \right) &= \psi_1 \left( \frac{ (a_1 \otimes b_1) \cdot (c_2 \otimes d_2)}{ (c_1 \otimes d_1) \cdot (c_2 \otimes d_2) } + \frac{ (a_2 \otimes b_2) \cdot (c_1 \otimes d_1)}{ (c_2 \otimes d_2) \cdot (c_1 \otimes d_1) } \right)\\
& = \psi_1 \left( \frac{1}{c_1 c_2 \otimes d_1 d_2} \cdot (a_1 c_2 \otimes b_1 d_2 + c_1 a_2 \otimes d_1 b_2) \right) \\
&= \left(\frac{1}{c_1 c_2} \otimes \frac{1}{d_1 d_2}\right) \cdot \left( a_1 c_2 \otimes b_1 d_2 + c_1 a_2 \otimes d_1 b_2 \right).
\end{align*}
	\end{proof}

\end{document}
